% Project: `` Anchor institutions as macroeconomic stabilizers"
% Section: Introduction

\section{ Introduction}

\noindent Universities are large, immobile institutions that shape the economic geography of the places in which they are embedded. They employ substantial numbers of workers, draw students who become temporary residents, and support a dense ecosystem of services, housing, and local suppliers. Classical regional economics views such institutions as anchors that generate local multipliers, stimulate non-tradable demand, and contribute to long-run development. At the same time, high dependence on a single anchor can create vulnerabilities, particularly when the anchor’s local footprint contracts. This tension motivates a broad question that has become increasingly relevant in urban and regional research: How do anchor institutions shape the spatial transmission of aggregate shocks, and under what conditions do they act as stabilizers or sources of local fragility?

\subsection{ Related Literature}
\noindent \textbf{Universities as anchor institutions.} A growing literature investigates how universities function as anchor institutions that shape local economic resilience, employment, and spatial development. Classical urban and regional studies view universities as immobile organizations that generate backward and forward linkages, local spending multipliers, and long-term knowledge spillovers \citep{Saxenian1996RegionalAdvantage,felsensteinUniversityMetropolitanArena1996}. Recent measurement advances make this idea operational. The Philadelphia Fed’s Anchor Economy Initiative documents which U.S. communities are most reliant on universities for employment and income, highlighting college towns where a single campus accounts for a large share of local GDP and jobs \citep{Harker2024AnchorReliance}. This reliance has important implications. University scale and in-person presence can transmit shocks to non-tradable sectors such as housing, retail, and restaurants, and can dampen downturns when university activity expands or remains stable. Overall, the literature distinguishes between the long-run structural impacts of universities on local development and their short-run cyclical role as stabilizers of regional economies.

\noindent \textbf{Universities and long-run local development.} In the long run, universities raise the local stock of skills and foster innovation, but the magnitude of these effects depends on graduate retention, industrial structure, and local absorption capacity. Evidence from the United States and Europe shows that while university openings or expansions attract students and increase educational attainment, they yield heterogeneous impacts on wages and employment, often concentrated in dynamic or high-tech regions \citep{berlingieriCollegeOpeningsLocal2022,ferhatImpactUniversityOpenings,amendolaDoesGraduateHuman2020,abelCollegesUniversitiesIncrease2012}. Other work highlights that research universities, through academic R\&D and knowledge spillovers, can modify the local occupational composition even when direct employment effects are modest \citep{beesonEffectsCollegesUniversities,harrisUniversitiesAnchorInstitutions2016}. These studies show that universities shape regional economies through human capital accumulation and persistent changes in industrial structure.

\noindent \textbf{The role of universities as cyclical stabilizers.} A second line of work focuses on universities’ short-run stabilizing role in local economies. Empirical studies in labor and education economics consistently find that schooling is countercyclical. When labor market prospects weaken, the opportunity cost of schooling declines and enrollments rise, especially among teens and prime-aged adults who invest in education to improve employability \citep{dellasBusinessCyclesSchooling2003,bettsSafePortStorm1995}. Because university payrolls and student spending inject relatively stable demand into local non-tradable sectors, this countercyclicality can act as a built-in stabilizer for college-town economies. In regional science, this mechanism aligns with the concept of resistance in economic resilience, defined as the capacity of a region to experience smaller initial declines in output or employment after a shock \citep{martinNotionRegionalEconomic2015,martinRegionalEconomicResilience2012}.

Consistent with this mechanism, recent papers provide direct empirical evidence. \citet{howardUniversitiesImproveLocal2024} show that historical variation in the location of regional public universities explains greater employment resilience during sectoral downturns, largely driven by consumption in local non-tradables, with faculty and student populations channeling external income into local spending. However, these effects are heterogeneous. Research-intensive campuses with substantial R\&D, medical centers, or graduate programs generate stronger and more durable local demand and spillovers, whereas smaller or teaching-focused colleges in thinner markets provide more limited stabilization \citep{howardUniversitiesImproveLocal2024,beesonEffectsCollegesUniversities}. Evidence from the university openings and expansions literature reinforces this heterogeneity, finding that instructional growth reliably increases student inflows and local human capital, but broader labor-market multipliers depend on local absorptive capacity and graduate retention \citep{berlingieriCollegeOpeningsLocal2022,ferhatImpactUniversityOpenings,amendolaDoesGraduateHuman2020}.

Conversely, the very dependence that dampens typical downturns may become a source of fragility when the anchor’s activity contracts. The regional resilience literature emphasizes that the nature of the shock matters. When disturbances directly suppress the anchor’s local footprint, such as temporary campus shutdowns or institutional closures, the stabilizing channel weakens and exposure to non-tradables becomes a liability \citep{martinNotionRegionalEconomic2015}. Case-based and emerging empirical work documents sharp contractions in services, retail activity, and rental markets when students abruptly depart \citep{howardUniversitiesImproveLocal2024,Harker2024AnchorReliance}. University towns therefore combine stability and vulnerability. They are cushioned when enrollment rises countercyclically or payrolls persist, yet vulnerable when the campus reduces or suspends in-person operations.

\noindent \textbf{Closest work to this study.} \citet{jumpResearchUniversitiesRecession} examine whether counties hosting state flagship universities experienced greater resilience across the 2001, 2008--2009, and 2020 U.S. recessions. Using synthetic difference in differences, they find no cushioning effect during the 2001 recession, a substantial stabilizing effect during the Great Recession, and a negative effect during 2020, when in-person activity collapsed. These findings are consistent with the stability and fragility mechanisms described above. This project is closely related in spirit, but it focuses on a broader set of college towns, incorporates measures of heterogeneity in university dependence, and examines labor market and housing outcomes across multiple geographic scales.

\subsection{ Motivation and Scope}

\noindent Universities offer a natural context for studying the role of anchor institutions in the stabilization of local economies or the amplification of downturns, because they combine stability through countercyclical enrollment and persistent payrolls with vulnerability when student presence or operational modes change. This project contributes to that broader inquiry by studying how university-dependent local economies respond to different types of national recessions, with a particular focus on whether cyclical stabilization disappears or reverses when in-person university activity becomes constrained.

The empirical work is still in development. At this stage, I compile unemployment data for counties and metropolitan areas from 1990 to 2024 and classify college towns using an external list from \citet{GumprechtBook}. Preliminary descriptive patterns suggest that college-town MSAs exhibit slightly lower unemployment rates and somewhat narrower cyclical swings, except during the COVID recession, when differences are less pronounced. Event-study regressions show small stabilizing effects for MSAs in the pre-COVID period and no clear differences at the county level. These results are preliminary and limited by the current college-town definition, which does not yet incorporate endogenous measures of university dependence.

Several components of the project remain under construction. I am in the process of assembling detailed QCEW industry-by-county employment data to build endogenous measures of university exposure based on higher education employment and payroll shares. I will merge these with IPEDS student and staffing data, LEHD commuting flows, ACS demographic controls, and housing market data from Zillow and the FHFA. The empirical strategy will combine synthetic difference in differences, interaction difference in differences, and local projections to estimate both average and dynamic effects of university dependence during recessions. A separate event-study will examine the 2020 campus shutdown in detail. Future stages will analyze sectoral employment, small-business sensitivity, and housing market adjustment, and will test mechanisms related to local service intensity and housing supply tightness.

The remainder of the paper is organized as follows. Section 2 describes the construction of the county- and MSA-level datasets, including measures of university exposure, local labor markets, and housing outcomes. Section 3 outlines the empirical strategies and identification assumptions. Section 4 presents the main results on unemployment, sectoral employment, and housing markets across the 2001, 2008--2009, and 2020 recessions, and explores heterogeneity by local structure. Section 5 concludes by discussing the implications for our understanding of anchor institutions and regional resilience, and by outlining directions for future work.











